% Formal π-Calculus Model for Agent Runtime Contracts
% For AAAI 2027 submission - Section 3

\documentclass{article}
\usepackage{amsmath,amssymb,amsthm}
\usepackage{stmaryrd}  % For semantic brackets
\usepackage{mathpartir} % For inference rules
\usepackage{xcolor}

\newtheorem{theorem}{Theorem}
\newtheorem{lemma}[theorem]{Lemma}
\newtheorem{definition}{Definition}

\title{Formal Semantics for Multi-Agent Contract-Based Coordination}
\author{Agent Runtime Team}
\date{\today}

\begin{document}
\maketitle

\section{Introduction}

We present a formal model for agent runtime contracts using the π-calculus, a process algebra for concurrent systems with mobile processes. Our model formalizes the composition of heterogeneous agent plugins and proves key safety properties (deadlock-freedom, progress guarantees).

\section{Syntax}

\subsection{Contract Channels}

Each contract type is modeled as a typed channel with input/output capabilities:

\begin{align*}
\text{Contract Types: } \tau ::=\quad 
  &\mathtt{ToolContract}(tool\_name, args, result) \\
  &| \mathtt{MessageContract}(target, msg, ack) \\
  &| \mathtt{SharedContext}(key, value) \\
  &| \mathtt{RecorderContract}(event) \\
  &| \mathtt{ControlContract}(signal)
\end{align*}

\subsection{Process Terms}

Agents and contracts are modeled as concurrent processes:

\begin{align*}
P, Q ::=\quad 
  &\mathbf{0} && \text{(nil process)} \\
  &| c!\langle v \rangle.P && \text{(output on channel $c$)} \\
  &| c?(x).P && \text{(input on channel $c$, bind to $x$)} \\
  &| P \mid Q && \text{(parallel composition)} \\
  &| \nu c.P && \text{(channel restriction)} \\
  &| !P && \text{(replication)} \\
  &| \tau.P && \text{(internal action)}
\end{align*}

\subsection{Agent Model}

An agent is a process composed of multiple contract channels:

\begin{align*}
\text{Agent}_i = \nu c_1, \ldots, c_n. (&\text{ToolContract}(c_1) \\
                                        &\mid \text{MessageContract}(c_2) \\
                                        &\mid \text{SharedContext}(c_3))
\end{align*}

\subsection{Contract Implementations}

Each contract is modeled as a replicated server process:

\paragraph{ToolContract}
\begin{align*}
\text{ToolContract}(c) = !c?(tool, args, ret). \tau. ret!\langle \text{execute}(tool, args) \rangle
\end{align*}

\paragraph{MessageContract}
\begin{align*}
\text{MessageContract}(c) = !c?(target, msg, ack). \tau. ack!\langle \text{sent} \rangle
\end{align*}

\paragraph{SharedContext}
\begin{align*}
\text{SharedContext}(c) &= !\text{Store}(store) \\
\text{Store}(store) &= c\_get?(key, ret). ret!\langle store[key] \rangle. \text{Store}(store) \\
                     &\quad + c\_set?(key, val). \text{Store}(store \cup \{key \mapsto val\})
\end{align*}

\section{Operational Semantics}

\subsection{Structural Congruence}

\begin{mathpar}
\inferrule{ }{P \mid \mathbf{0} \equiv P} \and
\inferrule{ }{P \mid Q \equiv Q \mid P} \and
\inferrule{ }{(P \mid Q) \mid R \equiv P \mid (Q \mid R)} \and
\inferrule{ }{\nu c.\mathbf{0} \equiv \mathbf{0}} \and
\inferrule{c \notin fn(P)}{\nu c.(P \mid Q) \equiv P \mid \nu c.Q}
\end{mathpar}

\subsection{Reduction Rules}

\begin{mathpar}
\inferrule[\textsc{Comm}]
  { }
  {c!\langle v \rangle.P \mid c?(x).Q \to P \mid Q\{v/x\}}
\and
\inferrule[\textsc{Par}]
  {P \to P'}
  {P \mid Q \to P' \mid Q}
\and
\inferrule[\textsc{Res}]
  {P \to P'}
  {\nu c.P \to \nu c.P'}
\and
\inferrule[\textsc{Struct}]
  {P \equiv P' \quad P' \to Q' \quad Q' \equiv Q}
  {P \to Q}
\end{mathpar}

\section{Contract Composition}

\subsection{Multi-Agent System}

A multi-agent system is the parallel composition of agents with shared contracts:

\begin{align*}
\text{MAS} = \nu shared. (&\text{Agent}_1(shared) \\
                          &\mid \text{Agent}_2(shared) \\
                          &\mid \text{SharedContext}(shared))
\end{align*}

\subsection{Communication Pattern}

Agent communication via MessageContract:

\begin{align*}
\text{Agent}_A &= msg_c!\langle B, \text{"task"} \rangle.ack?(resp). \text{Continue} \\
\text{Agent}_B &= msg_c?(sender, task).process(task).ack!\langle \text{"done"} \rangle
\end{align*}

\section{Type System}

\subsection{Channel Types}

Channels are typed to ensure correct usage:

\begin{align*}
\Gamma ::= \emptyset \mid \Gamma, c:T
\end{align*}

where $T$ is a channel type:
\begin{align*}
T ::=\quad 
  &\uparrow\!\tau && \text{(output capability)} \\
  &| \downarrow\!\tau && \text{(input capability)} \\
  &| \uparrow\!\tau \otimes \downarrow\!\tau && \text{(bidirectional)}
\end{align*}

\subsection{Typing Rules}

\begin{mathpar}
\inferrule[\textsc{T-Out}]
  {\Gamma \vdash v:\tau \quad \Gamma \vdash P}
  {\Gamma, c:\uparrow\!\tau \vdash c!\langle v \rangle.P}
\and
\inferrule[\textsc{T-In}]
  {\Gamma, x:\tau \vdash P}
  {\Gamma, c:\downarrow\!\tau \vdash c?(x).P}
\and
\inferrule[\textsc{T-Par}]
  {\Gamma_1 \vdash P \quad \Gamma_2 \vdash Q}
  {\Gamma_1 \circ \Gamma_2 \vdash P \mid Q}
\end{mathpar}

\section{Safety Properties}

\subsection{Theorem 1: Deadlock-Freedom}

\begin{theorem}[Deadlock-Freedom]
Let $\text{MAS} = \text{Agent}_1 \mid \cdots \mid \text{Agent}_n \mid \text{Contracts}$ be a well-typed multi-agent system. Then $\text{MAS}$ is deadlock-free, i.e., if $\text{MAS} \not\to$ then $\text{MAS} \equiv \mathbf{0}$.
\end{theorem}

\begin{proof}[Proof Sketch]
We prove by induction on the structure of $\text{MAS}$:

\paragraph{Base Case:} If $\text{MAS} = \mathbf{0}$, then trivially deadlock-free.

\paragraph{Inductive Case:} Assume $\text{MAS} = P \mid Q$ where $P, Q$ are deadlock-free.

\begin{enumerate}
\item \textbf{Contract Channels are Acyclic}: 
  By construction, contract dependencies form a DAG:
  \begin{itemize}
  \item Agents call contracts (agent $\to$ contract)
  \item Contracts never call other contracts directly
  \item No cyclic dependencies: agent $\not\to$ contract $\not\to$ agent
  \end{itemize}

\item \textbf{All Contract Calls are Replicated Servers}:
  Each contract is defined as $!c?(x).P$, meaning it can always receive messages. Therefore, no agent can block waiting for a contract.

\item \textbf{SharedContext is Lock-Free in Model}:
  In our π-calculus model, SharedContext operations are atomic messages (no explicit locks). The implementation uses asyncio.Lock but the abstract model is lock-free.

\item \textbf{MessageContract is Asynchronous}:
  Message sending does not block (fire-and-forget). Even with acknowledgments, the ack channel is provided by caller, so no circular wait.
\end{enumerate}

Therefore, $\text{MAS}$ cannot deadlock unless agents create cyclic dependencies, which our type system prevents via acyclicity constraint.
\end{proof}

\subsection{Theorem 2: Progress}

\begin{theorem}[Progress]
Let $\text{MAS}$ be a well-typed multi-agent system. If an agent calls a contract method, the contract responds within a finite number of reduction steps (assuming non-diverging contract implementations).
\end{theorem}

\begin{proof}[Proof Sketch]
By induction on the structure of contract calls:

\paragraph{ToolContract:}
\begin{align*}
\text{Agent} &= tool_c!\langle \text{"api"}, args, ret \rangle.ret?(result).P \\
\text{ToolContract} &= !tool_c?(tool, args, ret').\tau.ret'!\langle r \rangle
\end{align*}

Reduction sequence:
\begin{align*}
&\text{Agent} \mid \text{ToolContract} \\
&\to ret?(result).P \mid \tau.ret!\langle r \rangle && \text{(by Comm)} \\
&\to ret?(result).P \mid ret!\langle r \rangle && \text{(by internal action)} \\
&\to P\{r/result\} && \text{(by Comm)}
\end{align*}

Progress guaranteed in 3 steps (assuming $\tau$ terminates).

\paragraph{SharedContext:}
Get operation:
\begin{align*}
&c\_get!\langle key, ret \rangle.ret?(val).P \mid \text{Store}(store) \\
&\to ret?(val).P \mid ret!\langle store[key] \rangle.\text{Store}(store) \\
&\to P\{store[key]/val\} \mid \text{Store}(store)
\end{align*}

Progress in 2 steps.

Therefore, all contract calls make progress within finite steps.
\end{proof}

\subsection{Theorem 3: Sequential Consistency for SharedContext}

\begin{theorem}[Sequential Consistency]
SharedContext operations appear to execute atomically in some sequential order consistent with the program order of each agent.
\end{theorem}

\begin{proof}[Proof Sketch]
Our SharedContext model is:
\begin{align*}
\text{Store}(store) = c\_get?(k, ret).ret!\langle store[k] \rangle.\text{Store}(store) \\
                     + c\_set?(k, v).\text{Store}(store \cup \{k \mapsto v\})
\end{align*}

\paragraph{Key Observation}: The store is passed as a functional argument, creating a sequential chain:
\begin{align*}
\text{Store}(s_0) \to \text{Store}(s_1) \to \text{Store}(s_2) \to \cdots
\end{align*}

Each state transition is atomic (single reduction step), and the order is determined by the reduction strategy (e.g., left-to-right).

\paragraph{Implementation Note}: The Python implementation uses asyncio.Lock, which provides mutual exclusion. The π-calculus model abstracts this as atomic state transitions.

Therefore, SharedContext maintains sequential consistency.
\end{proof}

\section{Limitations of Current Model}

\subsection{What We Model}

\begin{itemize}
\item Contract interfaces (channels with typed messages)
\item Synchronous request-response patterns
\item Atomic SharedContext operations
\item Deadlock-freedom via acyclic dependencies
\end{itemize}

\subsection{What We Don't Model (Yet)}

\begin{itemize}
\item \textbf{Timeouts}: Need timed π-calculus ($\pi_t$-calculus)
\item \textbf{Failures}: Need failure calculus (crash/Byzantine)
\item \textbf{Distributed Coordination}: Need distributed π-calculus with locations
\item \textbf{Priority Queues}: MessageContract has priorities, not modeled here
\item \textbf{TTL Expiry}: SharedContext TTL requires timed semantics
\end{itemize}

\subsection{Future Extensions}

\paragraph{Timed π-Calculus:} Add time annotations to model timeouts and TTL:
\begin{align*}
P ::= \cdots \mid \text{delay}(t).P \mid \text{timeout}(t, P, Q)
\end{align*}

\paragraph{Probabilistic π-Calculus:} Model probabilistic contract failures:
\begin{align*}
P ::= \cdots \mid P \oplus_p Q && \text{(choose $P$ with prob $p$)}
\end{align*}

\paragraph{Spatial π-Calculus:} Model distributed agents with locations:
\begin{align*}
P ::= \cdots \mid loc[P] && \text{(process at location)}
\end{align*}

\section{Conclusion}

We have formalized agent runtime contracts using π-calculus and proved:
\begin{enumerate}
\item \textbf{Deadlock-Freedom}: Acyclic contract dependencies prevent circular wait
\item \textbf{Progress}: All contract calls complete in finite steps
\item \textbf{Sequential Consistency}: SharedContext operations are atomic
\end{enumerate}

These guarantees ensure that contract-based multi-agent systems are safe and correct by construction.

\paragraph{Next Steps:}
\begin{itemize}
\item Formalize timeouts and TTL using timed π-calculus
\item Model distributed coordination with location calculus
\item Implement TLA+ specifications for model checking
\item Extend to probabilistic contracts (failure handling)
\end{itemize}

\end{document}
